% ===== PRÉAMBULE CANONIQUE — à coller VERBATIM en tête de CHAQUE .tex =====
\documentclass[11pt,a4paper]{article}
\usepackage[utf8]{inputenc}\usepackage[T1]{fontenc}\usepackage{lmodern}
\usepackage[french]{babel}
\usepackage{amsmath,amssymb,mathtools}\usepackage{icomma}
\usepackage[version=4]{mhchem}          % \ce{...} : équations & formules
\usepackage{siunitx}\sisetup{locale=FR} % \qty{}{}, \si{}, \ang{}
\usepackage{chemfig}                    % \chemfig{...} : structures développées
\usepackage{xcolor}\usepackage{colortbl}
\usepackage{tikz}\usepackage{pgfplots}\pgfplotsset{compat=1.18}
\usepackage[most]{tcolorbox}\usepackage{enumitem}\usepackage{array,booktabs}
\usepackage[a4paper,margin=2.2cm]{geometry}
\usetikzlibrary{arrows.meta,calc,angles,quotes,positioning,patterns,backgrounds,shapes.geometric}
\definecolor{primary}{RGB}{222,49,99}\definecolor{primarydark}{RGB}{150,22,55}
\definecolor{defcol}{RGB}{0,114,178}\definecolor{methcol}{RGB}{52,92,148}
\definecolor{excol}{RGB}{0,135,90}\definecolor{attncol}{RGB}{200,40,55}
\tcbset{boxbase/.style={enhanced,breakable,boxrule=0.9pt,arc=2.5pt,
  left=3mm,right=3mm,top=2.3mm,bottom=2.3mm,fonttitle=\bfseries,coltitle=white}}
% --- boîtes TUTORIEL (noms EXACTS, ne pas renommer) ---
\newtcolorbox{cle}[1][]{boxbase,colback=defcol!6,colframe=defcol,
  title={\textsc{À retenir}\ifx\relax#1\relax\relax\else\textmd{ : \itshape #1}\fi}}
\newtcolorbox{methode}[1][]{boxbase,colback=methcol!7,colframe=methcol,sharp corners,
  title={\textsc{Méthode}\ifx\relax#1\relax\relax\else\textmd{ --- \itshape #1}\fi}}
\newtcolorbox{exemplebox}[1][]{boxbase,colback=excol!5,colframe=excol,
  title={\textsc{Exemple}\ifx\relax#1\relax\relax\else\textmd{ : \itshape #1}\fi}}
\newtcolorbox{piege}[1][]{boxbase,colback=attncol!6,colframe=attncol,
  title={\textsc{Piège}\ifx\relax#1\relax\relax\else\textmd{ : \itshape #1}\fi}}
\newtcolorbox{remarque}[1][]{boxbase,colback=black!4,colframe=black!45,
  title={\textsc{Remarque}\ifx\relax#1\relax\relax\else\textmd{ : \itshape #1}\fi}}
% --- boîtes EXERCICE / CORRIGÉ (noms EXACTS, auto-comptées) ---
\newtcolorbox[auto counter]{exo}[1][]{boxbase,colback=primary!4,colframe=primary,
  title={\textsc{Exercice}~\thetcbcounter\ifx\relax#1\relax\relax\else\textmd{ --- \itshape #1}\fi}}
\newtcolorbox[auto counter]{corrige}[1][]{boxbase,colback=excol!4,colframe=excol,
  title={\textsc{Corrigé}~\thetcbcounter\ifx\relax#1\relax\relax\else\textmd{ --- \itshape #1}\fi}}
\newcommand{\R}{\mathbb{R}}\newcommand{\N}{\mathbb{N}}\newcommand{\Z}{\mathbb{Z}}\newcommand{\ds}{\displaystyle}
% =========================================================================
\begin{document}
\sisetup{per-mode=symbol}
\begin{corrige}[écrire l'expression de $K_{\mathrm{ps}}$]
Chaque concentration ionique est élevée à son coefficient stœchiométrique ; le solide n'apparaît pas.
\begin{enumerate}[label=\alph*)]
  \item $K_{\mathrm{ps}} = [\ce{Ag+}]\,[\ce{Cl-}]$.
  \item $K_{\mathrm{ps}} = [\ce{Pb^{2+}}]\,[\ce{Br-}]^{2}$.
  \item $K_{\mathrm{ps}} = [\ce{Fe^{3+}}]\,[\ce{OH-}]^{3}$.
  \item $K_{\mathrm{ps}} = [\ce{Ca^{2+}}]\,[\ce{CO3^{2-}}]$.
  \item $K_{\mathrm{ps}} = [\ce{Ag+}]^{2}\,[\ce{CrO4^{2-}}]$.
\end{enumerate}
\end{corrige}

\begin{corrige}[coefficient ou exposant ?]
\begin{enumerate}[label=\alph*)]
  \item $\ce{Ag3PO4(s) <=> 3 Ag^{+}(aq) + PO4^{3-}(aq)}$.
  \item $K_{\mathrm{ps}} = [\ce{Ag+}]^{3}\,[\ce{PO4^{3-}}]$.
  \item Le coefficient $3$ de \ce{Ag+} dans l'équation devient un \textbf{exposant} dans
        $K_{\mathrm{ps}}$ (loi d'action de masse), \emph{pas} un facteur multiplicatif. On écrit donc
        $[\ce{Ag+}]^{3}$ et non $3\,[\ce{Ag+}]$.
\end{enumerate}
\end{corrige}

\begin{corrige}[vrai ou faux : la nature de $K_{\mathrm{ps}}$]
\begin{enumerate}[label=\alph*)]
  \item \textbf{Faux.} $K_{\mathrm{ps}}$ est une constante d'équilibre : il est \emph{sans unité}.
  \item \textbf{Vrai.} C'est la constante d'équilibre de la réaction de dissolution.
  \item \textbf{Vrai.} La solubilité est une concentration molaire, en \si{\mole\per\litre}.
  \item \textbf{Vrai.} Comme toute constante d'équilibre, $K_{\mathrm{ps}}$ varie avec la température
        (le plus souvent il augmente à chaud).
\end{enumerate}
\end{corrige}

\begin{corrige}[calculer $K_{\mathrm{ps}}$ à partir des concentrations]
Pour un sel $1{:}1$, $K_{\mathrm{ps}} = [\text{cation}][\text{anion}]$ = produit des deux
concentrations mesurées.
\begin{enumerate}[label=\alph*)]
  \item $K_{\mathrm{ps}} = (\num{1.3e-5})^{2} = \num{1.7e-10}$.
  \item $K_{\mathrm{ps}} = (\num{1.0e-5})^{2} = \num{1.0e-10}$.
  \item $K_{\mathrm{ps}} = (\num{1.0e-13})^{2} = \num{1.0e-26}$.
\end{enumerate}
\end{corrige}

\begin{corrige}[cas $MX_n$ : l'exposant est $n$]
$K_{\mathrm{ps}}(\mathrm{MX}_n) = [\ce{M^{n+}}]\,[\ce{X-}]^{n}$.
\begin{enumerate}[label=\alph*)]
  \item $K_{\mathrm{ps}} = [\ce{Ag+}]\,[\ce{Br-}]$.
  \item $K_{\mathrm{ps}} = [\ce{Ca^{2+}}]\,[\ce{F-}]^{2}$.
  \item $K_{\mathrm{ps}} = [\ce{Al^{3+}}]\,[\ce{OH-}]^{3}$.
\end{enumerate}
\end{corrige}
\end{document}
